\documentclass[11pt]{article}
\usepackage[margin=1in]{geometry}
\usepackage{amsmath,amssymb,amsthm}
\usepackage{hyperref}

\newtheorem{theorem}{Theorem}
\newtheorem{lemma}[theorem]{Lemma}
\newtheorem{proposition}[theorem]{Proposition}
\newtheorem{corollary}[theorem]{Corollary}
\newtheorem{conjecture}{Conjecture}

\title{Proof of Easlick's Conjecture on Minimal-Advantage Non-Transitive $n$-Sided Dice}
\author{Salem Levinstein}
\date{November 25, 2026}

\begin{document}
\maketitle

\begin{abstract}
Let $G_n$ be the directed graph whose vertices are all $\binom{2n-1}{n}$ non-standard $n$-sided dice with face values in $\{1,\dots,n\}$, with a directed edge $D_{\mathbf{x}} \to D_{\mathbf{y}}$ whenever $\mathbf{x}$ beats $\mathbf{y}$ by the smallest attainable winning probability strictly exceeding $1/2$, namely $\sum_{i,j=1}^n \mathbf{1}_{\{x_i > y_j\}} = \lfloor n^2/2 \rfloor + 1$. Easlick (\texttt{arXiv:2610.40321v2}, 2026) proved that $G_3,G_4,G_5$ are acyclic and that $G_6$ has maximum directed cycle length~$6$, and conjectured that $n=6$ is the largest $n \ge 3$ for which the maximum cycle length of $G_n$ equals~$n$. We prove this conjecture in full by constructing explicit simple directed cycles in $G_n$ of length $3(n-7)$ for every odd $n \ge 11$, length $3(n-10)$ for every even $n \ge 16$, and length $n+1$ for $n \in \{7,8,9,10,12,14\}$. Thus $G_n$ contains a simple directed cycle of length at least $n+1$ for every $n \ge 7$. The entire $n \ge 7$ proof (together with a $6$-cycle at $n=6$ and a generalized antiautomorphism for all $n$) is formally verified in Lean~4 using only the standard axioms.
\end{abstract}

\section{Introduction}

For an integer $n \ge 3$, a \emph{non-standard $n$-sided die} taking face values in $[n] := \{1,\dots,n\}$ is represented uniquely as a non-decreasing $n$-tuple $\mathbf{x} = (x_1,\dots,x_n) \in [n]^n$ with $1 \le x_1 \le \cdots \le x_n \le n$. Given two dice $\mathbf{x},\mathbf{y}$, let
\[
  W(\mathbf{x},\mathbf{y}) := \sum_{i=1}^n \sum_{j=1}^n \mathbf{1}_{\{x_i > y_j\}}
\]
be the number of face pairs on which $\mathbf{x}$ strictly beats $\mathbf{y}$. The smallest integer $w$ satisfying $w/n^2 > 1/2$ is
\[
  T(n) := \left\lfloor \frac{n^2}{2} \right\rfloor + 1 =
  \begin{cases}
    \frac{n^2}{2} + 1 & \text{if } n \text{ is even,} \\[3pt]
    \frac{n^2 + 1}{2} & \text{if } n \text{ is odd.}
  \end{cases}
\]
In~\cite{easlick2026}, Easlick defines the directed graph $G_n = (V_n, E_n)$ on the set $V_n$ of all $\binom{2n-1}{n}$ non-standard $n$-sided dice by placing a directed edge $\mathbf{x} \to \mathbf{y}$ if and only if $W(\mathbf{x},\mathbf{y}) = T(n)$. A simple directed cycle of length $L \ge 3$ in $G_n$ is a \emph{non-transitive cycle of length $L$ with uniform and minimal winning advantage}.

Easlick~\cite[Theorem~1]{easlick2026} proved by exhaustive enumeration that $G_3, G_4, G_5$ are acyclic and that $G_6$ has maximum cycle length~$6$ (attained by two unique $6$-cycles). After finding cycles of length $n+1$ for $7 \le n \le 12$, Easlick formulated the following conjecture.

\begin{conjecture}[{\cite[Conjecture~1]{easlick2026}}]
$n = 6$ is the largest $n \ge 3$ such that the set of non-standard $n$-sided dice admits a non-transitive cycle of maximum length $n$ with uniform and minimal winning advantage.
\end{conjecture}

In other words, letting $L_{\max}(G_n)$ denote the maximum length of a simple directed cycle in $G_n$, Conjecture~1 asserts that $L_{\max}(G_n) \ne n$ for all $n \ge 7$. We prove a stronger quantitative statement establishing $L_{\max}(G_n) \ge n+1$ for every $n \ge 7$.

\begin{theorem}\label{thm:main}
For every integer $n \ge 7$, the directed graph $G_n$ contains a simple directed cycle of length at least $n + 1$. Specifically:
\begin{enumerate}
  \item For every odd $n \ge 11$, $G_n$ contains a simple directed cycle of length $3(n - 7) \ge n + 1$.
  \item For every even $n \ge 16$, $G_n$ contains a simple directed cycle of length $3(n - 10) \ge n + 1$.
  \item For every $n \in \{7, 8, 9, 10, 12, 14\}$, $G_n$ contains a simple directed cycle of length $n + 1$.
\end{enumerate}
Consequently, Conjecture~1 of~\cite{easlick2026} holds, and $n = 6$ is the unique integer $n \ge 3$ satisfying $L_{\max}(G_n) = n$.
\end{theorem}

We also note that the involutive antiautomorphism exhibited by Easlick for $n = 6$~\cite[Prollary~1]{easlick2026} extends verbatim to every $n \ge 3$: the face-reversal map $f(x_1,\dots,x_n) = (n+1-x_n,\dots,n+1-x_1)$ satisfies $f(f(\mathbf{x})) = \mathbf{x}$ and $W(f(\mathbf{y}), f(\mathbf{x})) = W(\mathbf{x}, \mathbf{y})$ for all $\mathbf{x},\mathbf{y} \in V_n$.

\section{The 3-Die Zipper Construction}

The key observation is that when three dice families each consist of $n - 1$ fixed base faces in $\{1,\dots,b-1\}$ plus one free top face in $\{b,\dots,n\}$, the win count between two families depends on the free top faces only through the single indicator $\mathbf{1}_{\{u > v\}}$.

\begin{lemma}[3-Die Zipper]\label{lem:zipper}
Let $n \ge 3$ and $2 \le b \le n - 2$, and set $m := n - b - 1 \ge 1$. Suppose $A, B, C : \{b,\dots,n\} \to V_n$ are maps such that:
\begin{enumerate}
  \item the images $A(\{b,\dots,n\})$, $B(\{b,\dots,n\})$, and $C(\{b,\dots,n\})$ are pairwise disjoint, and each map $A, B, C$ is injective;
  \item $A(x) \to B(y)$ in $G_n$ whenever $x > y$;
  \item $B(y) \to C(z)$ in $G_n$ whenever $y \le z$;
  \item $C(z) \to A(x)$ in $G_n$ whenever $z > x$.
\end{enumerate}
Then $G_n$ contains a simple directed cycle of length $3m = 3(n - b - 1)$.
\end{lemma}

\begin{proof}
Define $3m$ vertices $D_0, D_1, \dots, D_{3m-1} \in V_n$ by setting, for each $q \in \{0, 1, \dots, m - 1\}$,
\[
  D_{3q} :=
  \begin{cases}
    B(b) & \text{if } q = 0, \\
    B(b + m - q) & \text{if } 1 \le q < m,
  \end{cases}
  \qquad
  D_{3q+1} := C(b + m + 1 - q),
  \qquad
  D_{3q+2} := A(b + m - q).
\]
Since $b + m + 1 = n$, every argument lies in $\{b,\dots,n\}$. As $q$ ranges over $\{0,\dots,m-1\}$, the arguments of $B$ are $b$ (at $q = 0$) and $b+m-1, \dots, b+1$ (at $q = 1,\dots,m-1$), which are $m$ distinct integers; likewise the arguments of $C$ are $b+m+1,\dots,b+2$ and those of $A$ are $b+m,\dots,b+1$. By condition~(1), all $3m$ dice are pairwise distinct.

For every $q \in \{0,\dots,m-1\}$:
\begin{itemize}
  \item $D_{3q} \to D_{3q+1}$ holds by~(3), since $b \le b+m+1$ (for $q = 0$) and $b+m-q \le b+m+1-q$ (for $q \ge 1$).
  \item $D_{3q+1} \to D_{3q+2}$ holds by~(4), since $b+m+1-q > b+m-q$.
  \item If $q < m - 1$, $D_{3q+2} \to D_{3(q+1)}$ is $A(b+m-q) \to B(b+m-q-1)$, which holds by~(2).
  \item If $q = m - 1$, the closing edge $D_{3m-1} \to D_0$ is $A(b+1) \to B(b)$, which holds by~(2).
\end{itemize}
Since $m \ge 1$, this yields a simple directed cycle of length $3m \ge 3$ in $G_n$.
\end{proof}

\section{Proof of Theorem~\ref{thm:main}}

We use exponential notation $v^r$ for $r$ repeated faces of value $v$.

\subsection{Odd $n = 2k + 1 \ge 11$}

Let $n = 2k + 1$ with $k \ge 5$, so $T(n) = 2k^2 + 2k + 1$. Take $b = 6$ and define for $x, y, z \in \{6,\dots,2k+1\}$:
\[
  A(x) = (1^k,\, 5^k,\, x), \qquad
  B(y) = (2^{k-1},\, 3^1,\, 4^k,\, y), \qquad
  C(z) = (2^1,\, 3^{2k-1},\, z).
\]
Each tuple has length $2k+1$ and non-decreasing entries in $[2k+1]$. The images are pairwise disjoint since $1$ occurs only in $A(x)$ and $4$ occurs in $B(y)$ but not in $C(z)$; each map is injective via its last coordinate. Evaluating the pairwise win counts directly gives:
\begin{align*}
  W(A(x), B(y)) &= k(2k) + 2k + \mathbf{1}_{\{x > y\}} = 2k^2 + 2k + \mathbf{1}_{\{x > y\}}, \\
  W(B(y), C(z)) &= 1\cdot 1 + k(2k) + 2k + \mathbf{1}_{\{y > z\}} = 2k^2 + 2k + 1 + \mathbf{1}_{\{y > z\}}, \\
  W(C(z), A(x)) &= (2k)k + 2k + \mathbf{1}_{\{z > x\}} = 2k^2 + 2k + \mathbf{1}_{\{z > x\}}.
\end{align*}
Hence $A(x) \to B(y)$ for $x > y$, $B(y) \to C(z)$ for $y \le z$, and $C(z) \to A(x)$ for $z > x$. By Lemma~\ref{lem:zipper} with $b = 6$, $G_n$ contains a simple directed cycle of length $3(n - 7)$. For all odd $n \ge 11$, $3(n - 7) - (n + 1) = 2n - 22 \ge 0$.

\subsection{Even $n = 2k \ge 16$}

Let $n = 2k$ with $k \ge 8$, so $T(n) = 2k^2 + 1$. Take $b = 9$ and define for $x, y, z \in \{9,\dots,2k\}$:
\[
  A(x) = (1^{k-1},\, 5^1,\, 8^{k-1},\, x), \qquad
  B(y) = (2^{k-3},\, 3^2,\, 4^1,\, 7^{k-1},\, y), \qquad
  C(z) = (2^1,\, 3^{k-2},\, 6^k,\, z).
\]
Each tuple has length $2k$ and non-decreasing entries in $[2k]$. The images are pairwise disjoint since $1$ occurs only in $A(x)$ and $7$ occurs in $B(y)$ but not in $C(z)$; each map is injective via its last coordinate. Direct evaluation yields:
\begin{align*}
  W(A(x), B(y)) &= 1 \cdot k + (k-1)(2k-1) + (2k-1) + \mathbf{1}_{\{x > y\}} = 2k^2 + \mathbf{1}_{\{x > y\}}, \\
  W(B(y), C(z)) &= 2 \cdot 1 + 1 \cdot (k-1) + (k-1)(2k-1) + (2k-1) + \mathbf{1}_{\{y > z\}} = 2k^2 + 1 + \mathbf{1}_{\{y > z\}}, \\
  W(C(z), A(x)) &= 1 \cdot (k-1) + (k-2)(k-1) + k \cdot k + (2k-1) + \mathbf{1}_{\{z > x\}} = 2k^2 + \mathbf{1}_{\{z > x\}}.
\end{align*}
Hence $A(x) \to B(y)$ for $x > y$, $B(y) \to C(z)$ for $y \le z$, and $C(z) \to A(x)$ for $z > x$. By Lemma~\ref{lem:zipper} with $b = 9$, $G_n$ contains a simple directed cycle of length $3(n - 10)$. For all even $n \ge 16$, $3(n - 10) - (n + 1) = 2n - 31 \ge 1$.

\subsection{Remaining Small Cases $n \in \{7, 8, 9, 10, 12, 14\}$}

For $n = 14$ ($k = 7$, $T(14) = 99$), setting $b = 8$ and
\[
  A_{14}(x) = (1^6, 4^1, 7^6, x), \qquad
  B_{14}(y) = (2^3, 3^4, 6^6, y), \qquad
  C_{14}(z) = (2^2, 3^4, 5^7, z)
\]
satisfies Lemma~\ref{lem:zipper} (since $4 \cdot 2 + 6 \cdot 13 + 13 = 99$ for $B_{14} \to C_{14}$, and $1 \cdot 7 + 6 \cdot 13 + 13 = 13 \cdot 6 + 7 \cdot 1 + 13 = 98$ for $A_{14} \to B_{14}$ and $C_{14} \to A_{14}$), producing a simple directed cycle of length $3(14 - 9) = 15 = 14 + 1$ in $G_{14}$.

For $n \in \{7, 8, 9, 10, 12\}$, the explicit simple directed cycles of length $n + 1$ recorded in~\cite[Examplemmas~1--2]{easlick2026} are directly checked in both our Lean~4 formalization and Python verifier. Together with the odd ($n \ge 11$) and even ($n \ge 16$) families, every $n \ge 7$ satisfies $L_{\max}(G_n) \ge n + 1$. \qed

\section{Formal Verification and Reproducibility}

The $n \ge 7$ part of Theorem~\ref{thm:main} (together with \texttt{n6\_has\_cycle\_6} and the general antiautomorphism) is formally verified in Lean~4 (\texttt{v4.33.1}, Mathlib~\texttt{v4.33.1}) at \url{https://github.com/s66lem/conjectures} in the module \texttt{Conjectures.EaslickNontransitiveDice}:
\begin{itemize}
  \item \texttt{wins\_eq\_card\_pairs}, \texttt{minAdvTarget\_gt\_half}, and \texttt{minAdvTarget\_le\_of\_gt\_half} verify that the formal definitions match~\cite[Equations~(1)--(5)]{easlick2026}.
  \item \texttt{flipDie\_involutive} and \texttt{flipDie\_antiautomorphism} prove the involutive antiautomorphism property for all $n$.
  \item \texttt{odd\_cycle\_exists} ($n \ge 11$ odd, length $3(n-7)$), \texttt{even\_cycle\_exists} ($n \ge 16$ even, length $3(n-10)$), \texttt{cycle\_7}--\texttt{cycle\_14}, and \texttt{easlick\_conjecture1} establish Theorem~\ref{thm:main} using only the standard axioms (\texttt{propext}, \texttt{Classical.choice}, \texttt{Quot.sound}) with zero \texttt{sorry} and zero \texttt{native\_decide}.
\end{itemize}
An independent Python~3 standard-library verifier (\texttt{verify.py}) additionally re-verifies Easlick's Theorem~1 exhaustively for $n = 3, 4, 5, 6$ via Tarjan's strongly connected components algorithm and checks all cycles for $7 \le n \le 101$.

\begin{thebibliography}{9}
\bibitem{easlick2026}
T.~Easlick, \emph{Computer-assisted results for $n$-sided non-transitive dice taking face values in $[n]$}, \texttt{arXiv:2610.40321v2} [math.CO], November 2026.
\end{thebibliography}

\end{document}
